\section{$\lambda$計算}

\begin{definition}
$V$を可算集合とし$V$の元を変数と呼ぶことにする.\ $\Lambda$を次の条件を満たす最小の集合とする.
\begin{enumerate}
    \item $V \subset \Lambda$
    \item ${^{\forall}}M, N \in \Lambda, (M N) \in \Lambda$
    \item ${^{\forall}}x \in V,\ {^{\forall}}M \in \Lambda, (\lambda x.M) \in \Lambda$
\end{enumerate}
$\Lambda$の元を\textbf{$\lambda$項}と呼ぶ.
\hfill $\square$
\end{definition}

BNF記法を用いて$\lambda$項$e$は次のように定義できる.
\begin{center}
$e ::= x \mid \lambda x. e \mid e e$
\end{center}

\begin{remark}
$M,N,P \in \Lambda$に対して$((M N) P)$を$M N P$と表記し$(\lambda x.(\lambda y.M))$を$\lambda x. \lambda y.M$と表記する.
\hfill $\square$
\end{remark}

\begin{definition}
$M \in \Lambda$の自由変数全体の集合$\textrm{FV}(M)$を次のようにして帰納的に定める.
\begin{enumerate}
    \item $\textrm{FV}(x) = \{x\}$
    \item $\textrm{FV}(M N) = \textrm{FV}(M) \cup \textrm{FV}(N)$
    \item $\textrm{FV}(\lambda x.M) = \textrm{FV}(M) \setminus \{x\}$
\end{enumerate}
\hfill $\square$
\end{definition}